% Options for packages loaded elsewhere
\PassOptionsToPackage{unicode}{hyperref}
\PassOptionsToPackage{hyphens}{url}
\PassOptionsToPackage{dvipsnames,svgnames,x11names}{xcolor}
\documentclass[
]{article}
\usepackage{xcolor}
\usepackage[margin=28mm]{geometry}
\usepackage{amsmath,amssymb}
\setcounter{secnumdepth}{5}
\usepackage{iftex}
\ifPDFTeX
  \usepackage[T1]{fontenc}
  \usepackage[utf8]{inputenc}
  \usepackage{textcomp} % provide euro and other symbols
\else % if luatex or xetex
  \usepackage{unicode-math} % this also loads fontspec
  \defaultfontfeatures{Scale=MatchLowercase}
  \defaultfontfeatures[\rmfamily]{Ligatures=TeX,Scale=1}
\fi
\usepackage{lmodern}
\ifPDFTeX\else
  % xetex/luatex font selection
\fi
% Use upquote if available, for straight quotes in verbatim environments
\IfFileExists{upquote.sty}{\usepackage{upquote}}{}
\IfFileExists{microtype.sty}{% use microtype if available
  \usepackage[]{microtype}
  \UseMicrotypeSet[protrusion]{basicmath} % disable protrusion for tt fonts
}{}
\makeatletter
\@ifundefined{KOMAClassName}{% if non-KOMA class
  \IfFileExists{parskip.sty}{%
    \usepackage{parskip}
  }{% else
    \setlength{\parindent}{0pt}
    \setlength{\parskip}{6pt plus 2pt minus 1pt}}
}{% if KOMA class
  \KOMAoptions{parskip=half}}
\makeatother
\usepackage{longtable,booktabs,array}
\newcounter{none} % for unnumbered tables
\usepackage{calc} % for calculating minipage widths
% Correct order of tables after \paragraph or \subparagraph
\usepackage{etoolbox}
\makeatletter
\patchcmd\longtable{\par}{\if@noskipsec\mbox{}\fi\par}{}{}
\makeatother
% Allow footnotes in longtable head/foot
\IfFileExists{footnotehyper.sty}{\usepackage{footnotehyper}}{\usepackage{footnote}}
\makesavenoteenv{longtable}
% definitions for citeproc citations
\NewDocumentCommand\citeproctext{}{}
\NewDocumentCommand\citeproc{mm}{%
  \begingroup\def\citeproctext{#2}\cite{#1}\endgroup}
\makeatletter
 % allow citations to break across lines
 \let\@cite@ofmt\@firstofone
 % avoid brackets around text for \cite:
 \def\@biblabel#1{}
 \def\@cite#1#2{{#1\if@tempswa , #2\fi}}
\makeatother
\newlength{\cslhangindent}
\setlength{\cslhangindent}{1.5em}
\newlength{\csllabelwidth}
\setlength{\csllabelwidth}{3em}
\newenvironment{CSLReferences}[2] % #1 hanging-indent, #2 entry-spacing
 {\begin{list}{}{%
  \setlength{\itemindent}{0pt}
  \setlength{\leftmargin}{0pt}
  \setlength{\parsep}{0pt}
  % turn on hanging indent if param 1 is 1
  \ifodd #1
   \setlength{\leftmargin}{\cslhangindent}
   \setlength{\itemindent}{-1\cslhangindent}
  \fi
  % set entry spacing
  \setlength{\itemsep}{#2\baselineskip}}}
 {\end{list}}
\usepackage{calc}
\newcommand{\CSLBlock}[1]{\hfill\break\parbox[t]{\linewidth}{\strut\ignorespaces#1\strut}}
\newcommand{\CSLLeftMargin}[1]{\parbox[t]{\csllabelwidth}{\strut#1\strut}}
\newcommand{\CSLRightInline}[1]{\parbox[t]{\linewidth - \csllabelwidth}{\strut#1\strut}}
\newcommand{\CSLIndent}[1]{\hspace{\cslhangindent}#1}
\setlength{\emergencystretch}{3em} % prevent overfull lines
\providecommand{\tightlist}{%
  \setlength{\itemsep}{0pt}\setlength{\parskip}{0pt}}
\usepackage{booktabs}
\usepackage{longtable}
\usepackage{microtype}
\usepackage{xcolor}
\usepackage{mathtools}
\definecolor{MidnightBlue}{RGB}{25,55,95}
\usepackage{bookmark}
\IfFileExists{xurl.sty}{\usepackage{xurl}}{} % add URL line breaks if available
\urlstyle{same}
\hypersetup{
  pdftitle={Smooth-point certificates for Polydegree containments},
  pdfauthor={Anonymous},
  colorlinks=true,
  linkcolor={MidnightBlue},
  filecolor={Maroon},
  citecolor={Blue},
  urlcolor={MidnightBlue},
  pdfcreator={LaTeX via pandoc}}

\title{Smooth-point certificates for Polydegree containments}
\usepackage{etoolbox}
\makeatletter
\providecommand{\subtitle}[1]{% add subtitle to \maketitle
  \apptocmd{\@title}{\par {\large #1 \par}}{}{}
}
\makeatother
\subtitle{A Jacobian interpretation of the Lewis--Perry--Straub
determinant}
\author{Anonymous}
\date{9 August 2026}

\begin{document}
\maketitle

{
\setcounter{tocdepth}{1}
\tableofcontents
}
\section*{Abstract}\label{abstract}
\addcontentsline{toc}{section}{Abstract}

Lewis, Perry and Straub introduced integral coefficient polynomials
\(g_{n,e}\) and an auxiliary determinant \(a_{d,e}\) whose simultaneous
specialisation yields containments between strata of plane polynomial
automorphisms. We make explicit a differential structure hidden in their
definitions. The entries of the determinant are precisely negative
partial derivatives,

\[
\alpha_{i,j,e}=-\partial_{x_{i+1}}g_{j+1,e},
\]

so \(a_{d,e}\) is, up to the sign \((-1)^e\), the Jacobian determinant
of \(g_{d,e},\ldots,g_{d+e-1,e}\). Weighted Euler identities then give
an exact integral syzygy. In the chart \(x_1=1\), this converts the
original non-vanishing condition into a smooth-point condition for
\(g_{d,e}=\cdots=g_{d+e-2,e}=0\) together with avoidance of the next
coefficient hypersurface.

This reformulation yields a smaller finite-field certificate. A simple
common zero modulo a prime with non-zero next coefficient lifts to
characteristic zero; no condition excluding primes dividing \(d+e-1\) is
needed. We also prove a curve-theoretic bridge in which a reduced affine
zero divisor supplies the required Jacobian point. Exact computations
for \(e=3\) show that, for \(2\le d\le20\), the entire affine complete
intersection is reduced and avoids the next coefficient. We also prove
that every one-variable boundary factor has only simple positive roots,
while adjacent boundary coprimality is checked through \(d=200\). These
results motivate, but do not prove, a uniform all-\(d\) affine
conjecture. The paper separates proved identities, inherited theorems,
exact finite evidence and open claims throughout.

\textbf{Keywords:} plane polynomial automorphisms; Polydegree
Conjecture; Jacobian criterion; multivariate Hensel lifting; weighted
homogeneous polynomials; complete intersections; computer-assisted
mathematics

\section{Introduction}\label{introduction}

The group of polynomial automorphisms of the complex affine plane
carries two structures whose interaction is subtle. The Jung--van der
Kulk decomposition assigns a polydegree sequence to each automorphism,
while the group is also an ind-variety in which one can ask when one
polydegree stratum lies in the Zariski closure of another. Lewis, Perry
and Straub recast a family of these containment questions in terms of
explicit coefficient polynomials and proved a specialisation theorem:
suitable complex values of \(x_1,\ldots,x_e\) imply

\[
\mathcal G_{(d+e)}\subseteq
\overline{\mathcal G_{(d,e+1)}}.
\tag{1.1}
\]

Their formulation makes the problem algorithmically decidable for fixed
\((d,e)\) and produced computer-aided ranges for \(e=3,4,5\) (Lewis et
al. 2019). Perry's earlier thesis developed the same computational
direction for the case later indexed by Lewis--Perry--Straub as \(e=3\)
(Perry 2016). The coefficient polynomials also sit within the rigidity
and Strong Factorial framework of Edo and van den Essen (Edo and Essen
2014).

The determinant appearing in the Lewis--Perry--Straub criterion was
described there as presently unmotivated (Lewis et al. 2019). The
starting point of this paper is that it is not an additional object: it
is the Jacobian determinant of the coefficient map already under study.
Once this observation is combined with weighted homogeneity, the
original specialisation criterion becomes a local geometric question:

\begin{quote}
Find one smooth point of the partial coefficient complete intersection
that does not lie on the next coefficient hypersurface.
\end{quote}

That change of language has concrete consequences. It removes the full
\(e\times e\) determinant from a finite-field search, permits useful
primes previously rejected by a residue-level determinant test, and
joins two seemingly different proof routes: Hensel lifting from a finite
field and divisor separation on a curve.

\subsection{Contributions}\label{contributions}

The proved contributions are as follows.

\begin{enumerate}
\def\labelenumi{\arabic{enumi}.}
\item
  We give a formal-inverse proof that every displayed rational
  multinomial formula for \(g_{n,e}\) is an element of
  \(\mathbf Z[x_1,\ldots,x_e]\).
\item
  We prove the derivative and determinant identities

  \[
  \alpha_{i,j,e}=-\frac{\partial g_{j+1,e}}{\partial x_{i+1}},
  \qquad
  a_{d,e}=(-1)^e\det D(g_{d,e},\ldots,g_{d+e-1,e}).
  \]
\item
  We derive an exact weighted-Euler syzygy, not merely an equality after
  evaluation on a zero locus. It identifies the Lewis--Perry--Straub
  non-vanishing condition with the non-vanishing of an affine Jacobian
  minor and the next coefficient.
\item
  We prove a reduced finite-field certificate theorem. A modular common
  zero at which the affine minor and the next coefficient are non-zero
  suffices. The prime may divide \(d+e-1\): the original determinant
  then vanishes in the residue field but remains non-zero at the
  characteristic zero lift.
\item
  We prove a prime--gap--Jacobian bridge on a smooth
  complete-intersection curve, including the affine-support and
  perfect-ground-field hypotheses needed for coefficient-one
  prime-divisor applications.
\item
  For \(e=3\) we prove uniform squarefreeness of the one-variable
  boundary factors by a terminating hypergeometric factorisation and
  finite free convolution. We also report exact finite computations that
  test a stronger scheme-theoretic conjecture. For \(2\le d\le20\),
  every affine common point is smooth and avoids the next coefficient,
  and adjacent boundary coprimality is checked through \(d=200\).
  Neither bounded sweep is promoted to an all-\(d\) theorem.
\end{enumerate}

\subsection{Relation to prior results and novelty
boundary}\label{relation-to-prior-results-and-novelty-boundary}

Lewis--Perry--Straub Theorem 13 records the ranges

\[
e=3,\ 2\le d<50;\qquad
e=4,\ 2\le d<20;\qquad
e=5,\ 2\le d\le12,
\]

while their Theorem 14, attributed to Edo, proves (1.1) whenever \(d-1\)
is divisible by \(e\) (Lewis et al. 2019). In particular, for \(e=3\)
the class \(d\equiv1\pmod3\) is already an infinite family. The present
paper therefore does not claim the first infinite family, nor does it
claim an all-\(d\) result.

A bounded search of the primary paper, Perry's thesis, the Strong
Factorial paper and indexed citing literature did not locate the
derivative--Jacobian packaging developed here. That is provisional
negative evidence, not proof of priority. The safe claim is that we
\textbf{make explicit and exploit} a Jacobian structure that was not
found in the sources inspected through 9 August 2026. Authenticated
specialist-database searches, original-author contact, expert review and
journal peer review remain outstanding.

\subsection{Evidence classes}\label{evidence-classes}

Four evidence classes are used below.

{\def\LTcaptype{none} % do not increment counter
\begin{longtable}[]{@{}
  >{\raggedright\arraybackslash}p{(\linewidth - 2\tabcolsep) * \real{0.5000}}
  >{\raggedright\arraybackslash}p{(\linewidth - 2\tabcolsep) * \real{0.5000}}@{}}
\toprule\noalign{}
\begin{minipage}[b]{\linewidth}\raggedright
Label
\end{minipage} & \begin{minipage}[b]{\linewidth}\raggedright
Meaning
\end{minipage} \\
\midrule\noalign{}
\endhead
\bottomrule\noalign{}
\endlastfoot
Proved & A general argument is included in this manuscript. \\
Inherited & A result is used from a cited external source, chiefly
Lewis--Perry--Straub Theorem 4. \\
Exactly checked & A finite statement was evaluated with exact arithmetic
and retained receipts. \\
Open & A conjecture or priority question is explicitly unresolved. \\
\end{longtable}
}

Internal replay, code-path diversity and mutation controls improve
confidence in the exact computations. They are not independent
reproduction, formal verification, expert acceptance or peer review.

\section{Coefficient polynomials and formal
inversion}\label{coefficient-polynomials-and-formal-inversion}

Fix \(e\ge1\). For \(\mathbf a=(a_1,\ldots,a_e)\in\mathbf Z_{\ge0}^e\),
write

\[
|\mathbf a|=a_1+\cdots+a_e,
\qquad
\mathbf a\!\cdot\!\mathbf N=a_1+2a_2+\cdots+ea_e.
\]

For \(n\ge0\), define

\[
g_{n,e}=\frac1{n+1}
\sum_{\mathbf a\cdot\mathbf N=n}
(-1)^{|\mathbf a|}
\binom{n+|\mathbf a|}{\mathbf a,n}\mathbf x^{\mathbf a}.
\tag{2.1}
\]

For \(0\le i\le e-1\) and \(j\ge-1\), define

\[
\alpha_{i,j,e}=
\sum_{\mathbf a\cdot\mathbf N=j-i}
(-1)^{|\mathbf a|}
\binom{|\mathbf a|+j+2}{\mathbf a,j+2}\mathbf x^{\mathbf a},
\tag{2.2}
\]

with an empty sum when \(j-i<0\), and put

\[
a_{d,e}=\det(\alpha_{i,j,e})_{
0\le i\le e-1,\ d-1\le j\le d+e-2}.
\tag{2.3}
\]

These are the definitions used by Lewis, Perry and Straub (Lewis et al.
2019). The factor \(1/(n+1)\) in (2.1) makes integrality worth stating
before reduction modulo a prime.

\subsection*{Lemma 2.1 (integrality and inverse coefficients)
{[}Proved{]}}\label{lemma-2.1-integrality-and-inverse-coefficients-proved}
\addcontentsline{toc}{subsection}{Lemma 2.1 (integrality and inverse
coefficients) {[}Proved{]}}

Let \(I(w)\) be the unique compositional inverse of

\[
H(t)=t+x_1t^2+\cdots+x_et^{e+1}.
\]

Then

\[
I(w)\in\mathbf Z[x_1,\ldots,x_e][[w]],
\qquad
g_{n,e}=[w^{n+1}]I(w).
\tag{2.4}
\]

\subsubsection*{Proof}\label{proof}
\addcontentsline{toc}{subsubsection}{Proof}

Write \(I(w)=w+\sum_{m\ge2}c_mw^m\). In the identity

\[
I+x_1I^2+\cdots+x_eI^{e+1}=w,
\]

the coefficient of \(w^m\) is \(c_m\) plus an integer polynomial in
\(x_1,\ldots,x_e,c_2,\ldots,c_{m-1}\). Induction therefore gives
\(c_m\in\mathbf Z[x_1,\ldots,x_e]\).

Formal Lagrange inversion gives

\[
[w^{n+1}]I(w)=\frac1{n+1}[t^n]
(1+x_1t+\cdots+x_et^e)^{-(n+1)}.
\]

Expanding the negative power multinomially yields exactly (2.1). Hence
the apparently rational expression is an integral polynomial.
\(\square\)

The recursion in the first paragraph and the multinomial expression in
the second are algebraically different constructions of the same
coefficients. They are used as separate internal replay paths in the
evidence package.

\section{The auxiliary determinant is a
Jacobian}\label{the-auxiliary-determinant-is-a-jacobian}

\subsection*{Proposition 3.1 (derivative identity)
{[}Proved{]}}\label{proposition-3.1-derivative-identity-proved}
\addcontentsline{toc}{subsection}{Proposition 3.1 (derivative identity)
{[}Proved{]}}

For \(0\le i\le e-1\) and \(j\ge-1\),

\[
\boxed{\alpha_{i,j,e}=-\frac{\partial g_{j+1,e}}{\partial x_{i+1}}.}
\tag{3.1}
\]

\subsubsection*{Proof}\label{proof-1}
\addcontentsline{toc}{subsubsection}{Proof}

Only terms of \(g_{j+1,e}\) whose \(x_{i+1}\) exponent is positive
survive differentiation. Put \(\mathbf b=\mathbf a+\mathbf e_{i+1}\).
Then

\[
\mathbf a\cdot\mathbf N=j-i,
\qquad
|\mathbf b|=|\mathbf a|+1.
\]

After multiplication by \(b_{i+1}=a_{i+1}+1\), the coefficient of
\(\mathbf x^{\mathbf a}\) in the derivative is

\[
\begin{aligned}
&\frac1{j+2}(-1)^{|\mathbf a|+1}
\frac{(j+1+|\mathbf b|)!}
{b_1!\cdots b_e!(j+1)!}(a_{i+1}+1)\\
&\hspace{18mm}=(-1)^{|\mathbf a|+1}
\frac{(j+2+|\mathbf a|)!}
{a_1!\cdots a_e!(j+2)!}.
\end{aligned}
\]

This is the negative of the coefficient in (2.2). Summing over the
common index set proves (3.1). \(\square\)

For \(d\ge2\) and \(e\ge2\), let

\[
D_{d,e}=\left(
\frac{\partial g_{d+r,e}}{\partial x_s}
\right)_{0\le r\le e-1,\ 1\le s\le e}.
\]

\subsection*{Corollary 3.2 (full determinant identity)
{[}Proved{]}}\label{corollary-3.2-full-determinant-identity-proved}
\addcontentsline{toc}{subsection}{Corollary 3.2 (full determinant
identity) {[}Proved{]}}

\[
\boxed{a_{d,e}=(-1)^e\det D_{d,e}.}
\tag{3.2}
\]

\subsubsection*{Proof}\label{proof-2}
\addcontentsline{toc}{subsubsection}{Proof}

The matrix in (2.3) is \(-D_{d,e}^{\mathsf T}\) by Proposition 3.1. Its
determinant is therefore \((-1)^e\det D_{d,e}\). \(\square\)

Thus \(a_{d,e}\) tests whether the full coefficient map

\[
(x_1,\ldots,x_e)\longmapsto
(g_{d,e},\ldots,g_{d+e-1,e})
\]

is locally invertible. The specialisation theorem does not, however,
require all \(e\) coefficients to vanish. The next section extracts the
relevant minor on the partial zero locus.

\section{Weighted Euler syzygy and the smooth-point
criterion}\label{weighted-euler-syzygy-and-the-smooth-point-criterion}

Give \(x_s\) weight \(s\). Each \(g_{n,e}\) is weighted homogeneous of
weight \(n\), so

\[
\sum_{s=1}^e sx_s\frac{\partial g_{n,e}}{\partial x_s}=ng_{n,e}.
\tag{4.1}
\]

Define the affine minor

\[
J_{d,e}=\det\left(
\frac{\partial g_{d+r,e}}{\partial x_s}
\right)_{0\le r\le e-2,\ 2\le s\le e}.
\tag{4.2}
\]

For \(0\le r\le e-1\), let \(M_r\) be the determinant of the
\((e-1)\times(e-1)\) submatrix of \(D_{d,e}\) obtained by deleting row
\(r\) and the \(x_1\) column. In particular, \(M_{e-1}=J_{d,e}\).

\subsection*{Proposition 4.1 (exact weighted-Euler syzygy)
{[}Proved{]}}\label{proposition-4.1-exact-weighted-euler-syzygy-proved}
\addcontentsline{toc}{subsection}{Proposition 4.1 (exact weighted-Euler
syzygy) {[}Proved{]}}

In \(\mathbf Z[x_1,\ldots,x_e]\),

\[
\boxed{
x_1a_{d,e}+(d+e-1)g_{d+e-1,e}J_{d,e}
=\sum_{r=0}^{e-2}(-1)^{e+r}(d+r)g_{d+r,e}M_r.}
\tag{4.3}
\]

Consequently,

\[
x_1a_{d,e}\equiv
-(d+e-1)g_{d+e-1,e}J_{d,e}
\pmod{(g_{d,e},\ldots,g_{d+e-2,e})}.
\tag{4.4}
\]

\subsubsection*{Proof}\label{proof-3}
\addcontentsline{toc}{subsubsection}{Proof}

In \(D_{d,e}\) replace the first column by

\[
C=\sum_{s=1}^e sx_s(D_{d,e})_s.
\]

By multilinearity, the resulting determinant is \(x_1\det D_{d,e}\):
every summand with \(s\ge2\) repeats an existing column. Equation (4.1)
makes entry \(r\) of \(C\) equal to \((d+r)g_{d+r,e}\). Expanding down
the first column and using (3.2) gives

\[
x_1a_{d,e}=\sum_{r=0}^{e-1}
(-1)^{e+r}(d+r)g_{d+r,e}M_r.
\]

The last summand is \(-(d+e-1)g_{d+e-1,e}J_{d,e}\). Moving it to the
left proves (4.3), and (4.4) follows. \(\square\)

\subsection*{Corollary 4.2 (scheme-theoretic smooth-point reformulation)
{[}Proved{]}}\label{corollary-4.2-scheme-theoretic-smooth-point-reformulation-proved}
\addcontentsline{toc}{subsection}{Corollary 4.2 (scheme-theoretic
smooth-point reformulation) {[}Proved{]}}

Let \(K\) be a field and let \(P\in K^e\) satisfy

\[
g_{d,e}(P)=\cdots=g_{d+e-2,e}(P)=0.
\]

Assume \(x_1(P)\ne0\) and \(\operatorname{char}K\nmid d+e-1\). Then the
following are equivalent:

\begin{enumerate}
\def\labelenumi{\arabic{enumi}.}
\item
  \(g_{d+e-1,e}(P)\ne0\) and \(a_{d,e}(P)\ne0\);
\item
  \(g_{d+e-1,e}(P)\ne0\) and \(J_{d,e}(P)\ne0\);
\item
  \(P\) lies outside \(V(g_{d+e-1,e})\), and the closed subscheme

  \[
  \operatorname{Spec}K[x_1,\ldots,x_e]/
  (g_{d,e},\ldots,g_{d+e-2,e})
  \]

  is smooth of codimension \(e-1\) at \(P\).
\end{enumerate}

\subsubsection*{Proof}\label{proof-4}
\addcontentsline{toc}{subsubsection}{Proof}

On the common zero locus, (4.4) makes (1) and (2) equivalent under the
stated unit hypotheses. For each of the first \(e-1\) coefficient
polynomials, (4.1) expresses its \(x_1\) derivative as a linear
combination of its \(x_2,\ldots,x_e\) derivatives because
\(x_1(P)\ne0\). Hence the partial Jacobian has rank \(e-1\) exactly when
the square minor (4.2) is non-zero. The Jacobian criterion for a
subscheme cut out by \(e-1\) equations in affine \(e\)-space gives the
equivalence with (3). \(\square\)

The characteristic restriction belongs to this residue-field
equivalence. It will disappear from the
finite-field-to-characteristic-zero theorem.

\section{A smaller Hensel
certificate}\label{a-smaller-hensel-certificate}

Work in the chart \(x_1=1\) and put

\[
F_r(x_2,\ldots,x_e)=g_{d+r,e}(1,x_2,\ldots,x_e).
\tag{5.1}
\]

The system \(F_0=\cdots=F_{e-2}=0\) has \(e-1\) equations in \(e-1\)
variables, and its Jacobian determinant is
\(J_{d,e}(1,x_2,\ldots,x_e)\).

\subsection*{Theorem 5.1 (smooth finite-field certificate)
{[}Proved{]}}\label{theorem-5.1-smooth-finite-field-certificate-proved}
\addcontentsline{toc}{subsection}{Theorem 5.1 (smooth finite-field
certificate) {[}Proved{]}}

Let \(d\ge2\) and \(e\ge2\). Suppose that a prime \(p\) and a point
\(Q\in\mathbf F_p^{e-1}\) satisfy

\[
F_0(Q)=\cdots=F_{e-2}(Q)=0,
\qquad
F_{e-1}(Q)J_{d,e}(1,Q)\ne0.
\tag{5.2}
\]

Then there is a complex specialisation satisfying the hypotheses of
Lewis--Perry--Straub Theorem 4. Consequently,

\[
\mathcal G_{(d+e)}\subseteq
\overline{\mathcal G_{(d,e+1)}}.
\tag{5.3}
\]

\subsubsection*{Proof}\label{proof-5}
\addcontentsline{toc}{subsubsection}{Proof}

The non-zero Jacobian determinant makes \(Q\) a simple zero of the
square system \(F_0=\cdots=F_{e-2}=0\). Lemma 2.1 shows that all these
polynomials have integral coefficients. Multivariate Hensel lifting
therefore produces an exact lift \(\widetilde Q\in\mathbf Z_p^{e-1}\) in
the same residue class. The values \(F_{e-1}(\widetilde Q)\) and
\(J_{d,e}(1,\widetilde Q)\) remain \(p\)-adic units. This is precisely
the simple-root setting of multivariate Hensel lifting (Grochow 2025).

Consider the localisation

\[
R=\frac{
\mathbf Q[x_2,\ldots,x_e,
1/(F_{e-1}J_{d,e})]
}{(F_0,\ldots,F_{e-2})},
\tag{5.4}
\]

where \(J_{d,e}\) is evaluated at \(x_1=1\). Evaluation at
\(\widetilde Q\) defines a unital \(\mathbf Q\)-algebra homomorphism
\(R\to\mathbf Q_p\), so \(R\ne0\). The extension
\(\mathbf Q\subset\mathbf C\) is faithfully flat; hence
\(R\otimes_{\mathbf Q}\mathbf C\ne0\) (Stacks project authors 2026a).
This is a non-zero finite-type \(\mathbf C\)-algebra. A maximal ideal
and the weak Nullstellensatz give a complex point at which the first
\(e-1\) equations vanish and \(F_{e-1}J_{d,e}\) does not (Stacks project
authors 2026b). Corollary 4.2, now in characteristic zero, shows that
\(a_{d,e}\) is non-zero there. The inherited Lewis--Perry--Straub
Theorem 4 yields (5.3) (Lewis et al. 2019). \(\square\)

The proof does \textbf{not} assert an embedding
\(\mathbf Q_p\hookrightarrow
\mathbf C\). The bridge is the non-zero localised \(\mathbf Q\)-algebra
and faithful scalar extension.

\subsection*{Remark 5.2 (bad-characteristic
rescue)}\label{remark-5.2-bad-characteristic-rescue}
\addcontentsline{toc}{subsection}{Remark 5.2 (bad-characteristic
rescue)}

No assumption \(p\nmid d+e-1\) is needed in Theorem 5.1. If
\(p\mid d+e-1\), equation (4.4) forces the residue of \(a_{d,e}\) to
vanish at a finite-field common zero. At the exact \(p\)-adic lift,
however,

\[
a_{d,e}(1,\widetilde Q)=
-(d+e-1)F_{e-1}(\widetilde Q)
J_{d,e}(1,\widetilde Q)\ne0
\quad\text{in }\mathbf Q_p.
\tag{5.5}
\]

The integer \(d+e-1\) may be a nonunit in \(\mathbf Z_p\), but it is not
zero in the characteristic-zero field \(\mathbf Q_p\). Thus the old
modular \(a_{d,e}\)-unit test is strictly stronger than the certificate
actually needed for complex existence.

\subsection*{Example 5.3 (exact positive regression) {[}Exactly
checked{]}}\label{example-5.3-exact-positive-regression-exactly-checked}
\addcontentsline{toc}{subsection}{Example 5.3 (exact positive
regression) {[}Exactly checked{]}}

For \(e=3\), \(d=8\), \(p=5\) and \((x_1,x_2,x_3)=(1,3,0)\), exact
finite-field evaluation gives

\[
g_{8,3}=g_{9,3}=0,\qquad
g_{10,3}=2,\qquad J_{8,3}=3,\qquad a_{8,3}=0
\pmod5.
\]

The legacy determinant-unit certificate rejects this point. Theorem 5.1
accepts it, and (5.5) proves that the determinant at the exact
\(5\)-adic lift is non-zero. The fixture is retained as a positive
regression test; it is not merely an illustrative calculation.

\section{A curve-theoretic route to the same
point}\label{a-curve-theoretic-route-to-the-same-point}

The smooth-point criterion also clarifies what a divisor argument must
prove. The statement below deliberately distinguishes an algebraically
closed geometric theorem from descent over a perfect ground field.

\subsection*{Theorem 6.1 (prime--gap--Jacobian bridge)
{[}Proved{]}}\label{theorem-6.1-primegapjacobian-bridge-proved}
\addcontentsline{toc}{subsection}{Theorem 6.1 (prime--gap--Jacobian
bridge) {[}Proved{]}}

Work over \(\mathbf C\) in the chart \(x_1=1\). Suppose the first
\(e-2\) polynomials

\[
F_0,\ldots,F_{e-3}
\]

cut out a smooth, geometrically integral complete-intersection curve
\(C\subset\mathbf A^{e-1}\) \textbf{scheme-theoretically}. Let
\(\overline C\) be its smooth projective completion, and regard the
\(F_r\) as rational functions on \(\overline C\). Assume:

\begin{enumerate}
\def\labelenumi{\arabic{enumi}.}
\tightlist
\item
  \(F_{e-2}|_{\overline C}\) is non-zero, and its zero divisor is
  non-empty, reduced and supported inside the affine open \(C\);
\item
  \(F_{e-1}(Q)\ne0\) at every point
  \(Q\in\operatorname{Supp}(F_{e-2}|_{\overline C})_0\).
\end{enumerate}

Then every point \(Q\) of that zero divisor satisfies

\[
F_0(Q)=\cdots=F_{e-2}(Q)=0,
\qquad
F_{e-1}(Q)J_{d,e}(1,Q)\ne0.
\tag{6.1}
\]

Consequently, each such point supplies a Lewis--Perry--Straub
specialisation.

\subsubsection*{Proof}\label{proof-6}
\addcontentsline{toc}{subsubsection}{Proof}

The support condition makes \(Q\) an affine point at which all displayed
functions are regular. Smooth complete-intersection status says that
\(dF_0,\ldots,dF_{e-3}\) are independent and cut the tangent space down
to the one-dimensional space \(T_QC\). Reducedness of the zero divisor
means that \(F_{e-2}|_C\) has a simple zero, equivalently
\(dF_{e-2}|_{T_QC}\ne0\). Hence \(dF_0,\ldots,dF_{e-2}\) are
independent. Their square Jacobian determinant is \(J_{d,e}(1,Q)\), so
it is non-zero. The second hypothesis supplies the other factor in
(6.1), and Corollary 4.2 gives the original determinant condition.
\(\square\)

The requirement that \(C\) be the scheme-theoretic complete intersection
is load-bearing. Smoothness of a selected component would not exclude
another component or embedded structure from changing the Jacobian
statement. The affine-support condition is equally necessary: a zero at
the projective boundary does not define a point in the chart used by
Theorem 5.1.

\subsection*{Corollary 6.2 (perfect-ground-field descent)
{[}Proved{]}}\label{corollary-6.2-perfect-ground-field-descent-proved}
\addcontentsline{toc}{subsection}{Corollary 6.2 (perfect-ground-field
descent) {[}Proved{]}}

Let \(k\subset\mathbf C\) be a perfect field of characteristic zero, and
choose its algebraic closure \(\overline k\) inside \(\mathbf C\).
Suppose the scheme-theoretic complete intersection \(C\), its smooth
projective completion and the \(F_r\) are defined over \(k\). Assume the
zero divisor of \(F_{e-2}\) is one closed point with coefficient one, is
supported in the affine open \(C\), and is disjoint from the zero and
pole support of \(F_{e-1}\). Then, after base change to \(\overline k\),
every geometric point above that closed point satisfies Theorem 6.1.

\subsubsection*{Proof}\label{proof-7}
\addcontentsline{toc}{subsubsection}{Proof}

Over a perfect field, the residue extension of a closed point is
separable. A coefficient-one closed point therefore becomes a disjoint
reduced collection of geometric points after base change. Affine support
and disjointness from the zero and pole support of \(F_{e-1}\) are
preserved by base change, so Theorem 6.1 applies. \(\square\)

Over an algebraically closed field, an individual closed point has
degree one. A degree-\(h\) coefficient-one prime certificate must
therefore be formulated over its ground field and then base-changed;
phrasing it only over \(\mathbf C\) loses the arithmetic degree
information.

\section{\texorpdfstring{Exact \(e=3\) geometry and the uniform
conjecture}{Exact e=3 geometry and the uniform conjecture}}\label{exact-e3-geometry-and-the-uniform-conjecture}

Set

\[
G_n(X,Y)=g_{n,3}(1,X,Y),\qquad
I_d=(G_d,G_{d+1})\subset\mathbf Q[X,Y],
\]

and

\[
J_d=\det\frac{\partial(G_d,G_{d+1})}{\partial(X,Y)}.
\]

The following section reports exact finite calculations. It is not used
to prove the general identities above.

\subsection*{Proposition 7.1 (verified affine range) {[}Exactly
checked{]}}\label{proposition-7.1-verified-affine-range-exactly-checked}
\addcontentsline{toc}{subsection}{Proposition 7.1 (verified affine
range) {[}Exactly checked{]}}

For every \(2\le d\le20\), exact Gr"obner-basis computation over
\(\mathbf Q\) establishes:

\begin{enumerate}
\def\labelenumi{\arabic{enumi}.}
\item
  \(g_{d,3}\) and \(g_{d+1,3}\) have homogeneous greatest common divisor
  \(1\);
\item
  \(I_d\) is zero-dimensional and

  \[
  \dim_{\mathbf Q}\mathbf Q[X,Y]/I_d
  =\left\lfloor\frac{d(d+1)}6\right\rfloor;
  \tag{7.1}
  \]
\item
  \(I_d+(J_d)=(1)\);
\item
  \(I_d+(G_{d+2})=(1)\).
\end{enumerate}

The quotient lengths for \(d=2,\ldots,20\) are

\[
1,2,3,5,7,9,12,15,18,22,26,30,35,40,45,51,57,63,70.
\]

Because \(I_d\) is a height-two complete intersection in the smooth
surface \(\mathbf A^2\), the third identity says that every geometric
point has full Jacobian rank. The entire affine scheme is therefore
reduced, indeed finite \textquotesingle\{e\}tale over \(\mathbf Q\). The
fourth identity says scheme-theoretically that no affine common point
lies on \(G_{d+2}=0\). Thus every affine point in the verified range is
a valid smooth specialisation.

\subsection{Newton polygons and the observed
length}\label{newton-polygons-and-the-observed-length}

The Newton polygon of \(G_n\) is

\[
\Delta_n=\operatorname{conv}\{(b,c)\in\mathbf Z_{\ge0}^2:2b+3c\le n\}.
\]

An exact six-residue-class calculation gives

\[
\operatorname{MV}(\Delta_d,\Delta_{d+1})
=\left\lfloor\frac{d(d+1)}6\right\rfloor.
\tag{7.2}
\]

The observed quotient length therefore reaches the Newton mixed-area
count. This suggests a toric non-degeneracy proof, but the specific
coefficients of \(G_d,G_{d+1}\) must be controlled; generic support
geometry alone is not a proof for this pair.

\subsection{Boundary reduction}\label{boundary-reduction}

Put

\[
B_n(X,Y)=g_{n,3}(0,X,Y).
\]

Let \(\epsilon=n\bmod2\), \(A=(n-3\epsilon)/2\), and \(z=Y^2/X^3\). On
the boundary torus \(XY\ne0\),

\[
B_n=X^AY^\epsilon P_n(z),
\]

where

\[
P_n(z)=\frac1{n+1}
\sum_{k=0}^{\lfloor A/3\rfloor}
(-1)^{A+\epsilon-k}
\frac{(n+A+\epsilon-k)!}
{n!(A-3k)!(\epsilon+2k)!}z^k.
\tag{7.3}
\]

The adjacent coefficient ratio is

\[
\frac{[z^{k+1}]P_n}{[z^k]P_n}
=-
\frac{(A-3k)(A-3k-1)(A-3k-2)}
{(n+A+\epsilon-k)(\epsilon+2k+1)(\epsilon+2k+2)}.
\tag{7.4}
\]

The minimum exponents give a uniform elementary statement: the common
monomial factor of \(B_d\) and \(B_{d+1}\) is

\[
X\quad\text{if }d\equiv1\pmod3,
\qquad 1\quad\text{otherwise}.
\tag{7.5}
\]

No power \(X^2\) and no \(Y\) divides both. Ruling out a further torus
factor is separate.

\subsection*{Theorem 7.2 (uniform boundary squarefreeness)
{[}Proved{]}}\label{theorem-7.2-uniform-boundary-squarefreeness-proved}
\addcontentsline{toc}{subsection}{Theorem 7.2 (uniform boundary
squarefreeness) {[}Proved{]}}

For every \(n\ge2\), every zero of \(P_n\) is real, positive and simple.
In particular, \(P_n\) is squarefree over every characteristic-zero
field.

\subsubsection*{Proof}\label{proof-8}
\addcontentsline{toc}{subsubsection}{Proof}

Put \(M=\lfloor A/3\rfloor\), write \(A=3M+\delta\) with
\(\delta\in\{0,1,2\}\), and let \(c_{n,M}\) be the non-zero leading
coefficient of \(P_n\). Set

\[
N=n+A+\epsilon,\qquad H=N-M+1,
\qquad
\sigma=\begin{cases}1/2,&\epsilon=0,\\-1/2,&\epsilon=1,\end{cases}
\qquad a=-M+\sigma,
\]

and

\[
(\beta_1,\beta_2)=
\begin{cases}
(1/3,2/3),&\delta=0,\\
(2/3,4/3),&\delta=1,\\
(4/3,5/3),&\delta=2.
\end{cases}
\]

Exact division of consecutive coefficients in (7.3) gives the reciprocal
identity

\[
\frac{s^MP_n(1/s)}{c_{n,M}}
={}_3F_2\!\left(
\begin{matrix}-M,\,a,\,H\\ \beta_1,\,\beta_2\end{matrix};
-\frac{4s}{27}\right).
\tag{7.6}
\]

Let

\[
F_n(x)={}_3F_2\!\left(
\begin{matrix}-M,\,a,\,H\\ \beta_1,\,\beta_2\end{matrix};x\right).
\]

The hypergeometric concatenation theorem for degree-\(M\) multiplicative
finite free convolution gives, up to a non-zero scalar,

\[
F_n=p_n\boxtimes_Mq_n,
\quad
p_n={}_2F_1(-M,a;\beta_1;x),
\quad
q_n={}_2F_1(-M,H;\beta_2;x)
\tag{7.7}
\]

(Martínez-Finkelshtein et al. 2024b). The
\href{https://dlmf.nist.gov/15.8.E1}{Pfaff transformation (DLMF
15.8.1)}, the \href{https://dlmf.nist.gov/18.5.E7}{Jacobi representation
(DLMF 18.5.7)}, and the \href{https://dlmf.nist.gov/18.16.E1}{strict
ordering of Jacobi zeros (DLMF 18.16.1)} show that \(p_n\) has \(M\)
simple negative roots and \(q_n\) has \(M\) simple roots in \((0,1)\).
The relevant Jacobi parameter pairs are

\[
(\beta_1-1,-\sigma),
\qquad
(\beta_2-1,H-M-\beta_2),
\]

and every entry exceeds \(-1\) (Olver et al. 2026). Sign preservation
for finite free convolution therefore makes \(F_n\) negative-rooted.

It remains to exclude repeated roots. Reflect
\(\widehat p_n(x)=p_n(-x)\). For \(M\ge2\) its positive roots are
simple, so its logarithmic mesh is greater than one. Proposition 2.17 of
Martínez-Finkelshtein, Morales and Perales gives

\[
\operatorname{lmesh}(\widehat p_n\boxtimes_Mq_n)
\ge \operatorname{lmesh}(\widehat p_n)>1
\]

(Martínez-Finkelshtein et al. 2024a). Hence the reflected convolution
has simple positive roots; the reflection identity makes \(F_n\) simple
and negative-rooted. The cases \(M=0,1\) are respectively constant and
linear. Finally, the scale in (7.6) and reciprocation turn the roots
into simple positive roots of \(P_n\). \(\square\)

The logarithmic-mesh step is load-bearing. A broader strict-interlacing
wording for a merely non-negative-rooted convolution factor admits a
degenerate counterexample; that formulation is not used here.

\subsection*{Proposition 7.3 (verified adjacent boundary range)
{[}Exactly
checked{]}}\label{proposition-7.3-verified-adjacent-boundary-range-exactly-checked}
\addcontentsline{toc}{subsection}{Proposition 7.3 (verified adjacent
boundary range) {[}Exactly checked{]}}

For every \(2\le d\le200\),

\[
\gcd(P_d,P_{d+1})=1,
\]

so through this range the full boundary gcd is exactly the factor in
(7.5).

When \(d\equiv1\pmod3\), the boundary support is the coarse point
\([0:0:1]\). On the \(\mu_3\) uniformising cover of the weighted stack
\(\mathcal P(1,2,3)\), the lowest terms of the two equations are
independent linear coordinates. The intersection is transverse on that
cover and contributes \(1/3\) to the stack intersection number. This
explains the floor in (7.1); \(1/3\) is an orbifold intersection
contribution, not a fractional coarse-scheme length.

\subsection*{Conjecture 7.4 (uniform affine smoothness)
{[}Open{]}}\label{conjecture-7.4-uniform-affine-smoothness-open}
\addcontentsline{toc}{subsection}{Conjecture 7.4 (uniform affine
smoothness) {[}Open{]}}

For every \(d\ge2\), the ideal \((G_d,G_{d+1})\) is radical of length

\[
\left\lfloor\frac{d(d+1)}6\right\rfloor,
\]

and both \(J_d\) and \(G_{d+2}\) are units in its coordinate algebra.

This would prove more than existence: every affine common point would be
a valid smooth specialisation. The exact ranges above strongly constrain
a future proof, but extending those ranges alone would add little
mathematics. The open structural tasks are uniform adjacent boundary
coprimality, toric non-degeneracy or radicality, Jacobian non-vanishing,
and exclusion of a triple common zero.

\section{Computational assurance and
reproducibility}\label{computational-assurance-and-reproducibility}

\subsection{Frozen finite-field
corpus}\label{frozen-finite-field-corpus}

The upstream v0.3.0 archive contains one finite-field witness for each
\(50\le d\le100\). Its witness JSON is byte-identical to the earlier
v0.2.0 corpus audited in this project. A fresh reconciliation on 9
August 2026 established:

{\def\LTcaptype{none} % do not increment counter
\begin{longtable}[]{@{}lr@{}}
\toprule\noalign{}
Gate & Result \\
\midrule\noalign{}
\endhead
\bottomrule\noalign{}
\endlastfoot
v0.3.0 manifest & 38/38 files match \\
upstream dual-path verifier & 51/51 under ordinary Python \\
upstream dual-path verifier & 51/51 under \texttt{python\ -O} \\
independent standard-library symbolic gates & 10/10 in both modes \\
independent witness relations & 51/51 in both modes \\
independent reports & byte-identical across modes \\
\end{longtable}
}

The independent verifier does not import the archive code. It uses
sparse integer dictionaries, a generic permutation determinant and
separate weighted-composition enumeration. This is internal code-path
diversity, not reproduction by an unaffiliated group.

\subsection{Negative controls}\label{negative-controls}

Sixteen mutation or negative controls and the positive
bad-characteristic fixture of Example 5.3 were exercised under ordinary
and optimised Python. They cover wrong derivative and determinant signs,
shifted indices, ordinary in place of weighted Euler coefficients,
corrupt residues and points, composite moduli, missing rows, singular
points, malformed types and an injected non-integral coefficient. All
expected outcomes were observed in both modes.

\subsection{A retained fail-open
defect}\label{a-retained-fail-open-defect}

The upstream verifier still contains a proof-critical Python
\texttt{assert} guarding coefficient divisibility before floor division.
Optimised Python removes that guard. An injected non-integral numerator
is rejected in ordinary mode but silently floor-divided under
\texttt{python\ -O}. The unmodified witness corpus passes in both modes,
and the independent verifier uses explicit exceptions; there is no
evidence that a current row is false. The defect is nevertheless
retained in the record because it is a real failure-detection weakness
in the upstream release.

\subsection{Geometry receipts}\label{geometry-receipts}

The affine sweep uses exact arithmetic over \(\mathbf Q\) with two
independent constructions of each \(G_n\) and Singular for ideal
computations. The boundary sweep uses exact univariate arithmetic.
Ordinary and optimised Python runs produced byte-identical JSON
receipts. The evidence archive records software versions, commands and
SHA-256 hashes.

\subsection{Assurance boundary}\label{assurance-boundary}

The package establishes that the stated scripts replay the stated finite
relations and that the written identities survived adversarial sign and
index checks. It does not establish novelty, priority, independent
reproduction, proof-assistant verification, specialist acceptance,
journal acceptance or the open all-\(d\) conjecture.

\section{Discussion}\label{discussion}

\subsection{What the Jacobian reformulation
changes}\label{what-the-jacobian-reformulation-changes}

The determinant identity is elementary once the two definitions are
aligned, but its value is organisational. It compresses three layers:

\[
\text{auxiliary determinant}
\longrightarrow
\text{Jacobian rank}
\longrightarrow
\text{smooth point avoiding one hypersurface}.
\]

For computation, the full determinant and its alpha-matrix
implementation can be removed from the search criterion. For geometry,
one may use any method that produces a reduced point away from the next
coefficient. For arithmetic, a simple residue-class point can be lifted
even when the original determinant vanishes modulo the chosen prime.

The three residues \(g_{d+e-1,e}\), \(J_{d,e}\) and \(a_{d,e}\) may
still be evaluated redundantly to catch software errors. They are not
algebraically independent on the common-zero locus: equation (4.4)
couples them exactly.

\subsection{What remains difficult}\label{what-remains-difficult}

Internal reproducibility is no longer the main bottleneck. The decisive
open problem is structural: prove that at least one good affine point
exists for all \(d\) in an unresolved column, or prove the stronger
Conjecture 7.4. Three routes are now sharply separated:

\begin{enumerate}
\def\labelenumi{\arabic{enumi}.}
\tightlist
\item
  prove uniform boundary coprimality and toric non-degeneracy;
\item
  exclude triple common zeros and singular intersections by recurrence
  or resultant identities;
\item
  construct a simple point asymptotically and certify a finite
  complement.
\end{enumerate}

Failure of this particular smooth-point route at one \(d\) would not
disprove the Polydegree Conjecture. Lewis--Perry--Straub Theorem 4 gives
a sufficient condition (Lewis et al. 2019).

\subsection{Significance and priority}\label{significance-and-priority}

The paper's current contribution is a reusable theorem package: a
transparent Jacobian interpretation, an exact syzygy, a sharper Hensel
criterion and a curve-theoretic bridge. The finite \(e=3\) geometry
identifies a precise frontier rather than another arbitrary cutoff. A
uniform theorem would materially increase significance; a larger finite
table would not.

Priority remains provisional. No inspected source states the present
derivative--determinant--Euler chain, but the first identity is close to
the multinomial reindexing already used in Lewis--Perry--Straub, and
elementary observations are particularly vulnerable to unpublished or
implicit prior use. Original-author contact and specialist review are
therefore substantive next steps, not administrative decoration.

\section{Limitations}\label{limitations}

\begin{enumerate}
\def\labelenumi{\arabic{enumi}.}
\tightlist
\item
  The Jacobian identities and bridge theorems have written proofs and
  internal adversarial review, but no proof-assistant formalisation or
  external specialist validation.
\item
  The prior-art search is bounded. Authenticated MathSciNet, zbMATH and
  Scopus searches and author contact were not completed.
\item
  The full affine statement is checked only for \(2\le d\le20\), and
  adjacent boundary coprimality only through \(d=200\). Uniform
  squarefreeness of each individual boundary factor does not imply
  adjacent coprimality. The finite bounds are evidence boundaries, not
  mathematical thresholds.
\item
  The 51 modular witnesses and their reimplementations originate within
  one research workflow. Different code paths do not constitute
  unaffiliated reproduction.
\item
  The manuscript inherits the correctness and interpretation of
  Lewis--Perry--Straub Theorem 4. It reproduces the hypotheses used but
  does not reprove their polynomial-automorphism construction.
\item
  No claim is made that the work currently meets a particular REF grade.
  Originality, significance and rigour require separate assessment, and
  publication infrastructure does not determine any of them.
\end{enumerate}

\section{Conclusion}\label{conclusion}

The Lewis--Perry--Straub determinant is the Jacobian of their
coefficient map. Weighted homogeneity turns that identity into an exact
syzygy and a smooth-point reformulation of the specialisation criterion.
The resulting finite-field certificate is smaller and strictly more
permissive: it needs a simple common zero and a non-zero next
coefficient, even when the prime divides \(d+e-1\). The same formulation
permits a divisor on a smooth curve to produce the required point.

For \(e=3\), exact calculations show a notably strong finite
pattern---the entire affine intersection is good through
\(d=20\)---while the boundary factors are uniformly squarefree. These
results isolate a clean uniform conjecture. The responsible next advance
is a structural proof or counterexample, not a longer table. Until then,
the general Jacobian and Hensel theorems stand as the paper's proved
contribution, while the all-\(d\) statement remains open.

\section*{Declarations}\label{declarations}
\addcontentsline{toc}{section}{Declarations}

\subsection*{Data and code
availability}\label{data-and-code-availability}
\addcontentsline{toc}{subsection}{Data and code availability}

All scripts, frozen witness data, exact JSON receipts, claim-status
metadata and SHA-256 manifests used for this candidate are included in
the accompanying reproducibility package. No DOI or public repository is
claimed at this stage. Public URLs and immutable identifiers must be
added only after a separately verified release.

\subsection*{Ethics statement}\label{ethics-statement}
\addcontentsline{toc}{subsection}{Ethics statement}

This is theoretical and computer-assisted mathematics. It involved no
human participants, personal data, animals or clinical materials;
institutional ethics approval was not applicable.

\subsection*{Funding}\label{funding}
\addcontentsline{toc}{subsection}{Funding}

An accountable-author funding attestation has not yet been supplied. No
funding status is asserted in this pre-deposit candidate; the
declaration must be completed and approved by the accountable author
before deposit.

\subsection*{Competing interests}\label{competing-interests}
\addcontentsline{toc}{subsection}{Competing interests}

An accountable-author competing-interest attestation has not yet been
supplied. No absence-of-conflict claim is asserted in this pre-deposit
candidate; the declaration must be completed and approved before
deposit.

\subsection*{Author contributions
(CRediT)}\label{author-contributions-credit}
\addcontentsline{toc}{subsection}{Author contributions (CRediT)}

An accountable-author CRediT attestation has not yet been supplied. No
contributor-role allocation is asserted in this pre-deposit candidate;
the statement must be completed and approved before deposit.

\subsection*{AI and automated-tool
disclosure}\label{ai-and-automated-tool-disclosure}
\addcontentsline{toc}{subsection}{AI and automated-tool disclosure}

Automated research, symbolic-computation and language-model tools
assisted literature triage, proof exploration, exact verification,
adversarial review, drafting and package preparation under human
direction. They are not listed as authors. The human author is
responsible for the submitted claims, source selection, final text and
release decision. Agreement among AI-assisted checks is not represented
as independent human review.

\subsection*{Acknowledgements}\label{acknowledgements}
\addcontentsline{toc}{subsection}{Acknowledgements}

The candidate builds on the definitions and specialisation theorem of
Drew Lewis, Kaitlyn Perry and Armin Straub. No endorsement by those
authors is implied.

\section*{References}\label{references}
\addcontentsline{toc}{section}{References}

\protect\phantomsection\label{refs}
\begin{CSLReferences}{1}{1}
\bibitem[\citeproctext]{ref-edo2014}
Edo, Eric, and Arno van den Essen. 2014. {``The {Strong Factorial
Conjecture}.''} \emph{Journal of Algebra} 397: 443--56.
\url{https://doi.org/10.1016/j.jalgebra.2013.09.011}.

\bibitem[\citeproctext]{ref-grochow2025}
Grochow, Joshua A. 2025. {``Multivariate {Hensel} Lifting Without
Assuming as Many Variables as Equations.''} \emph{The American
Mathematical Monthly} 132 (4): 350--55.
\url{https://doi.org/10.1080/00029890.2024.2440297}.

\bibitem[\citeproctext]{ref-lps2019}
Lewis, Drew, Kaitlyn Perry, and Armin Straub. 2019. {``An Algorithmic
Approach to the {Polydegree Conjecture} for Plane Polynomial
Automorphisms.''} \emph{Journal of Pure and Applied Algebra} 223 (12):
5346--59. \url{https://doi.org/10.1016/j.jpaa.2019.04.002}.

\bibitem[\citeproctext]{ref-martinez2024real}
Martínez-Finkelshtein, Andrei, Rafael Morales, and Daniel Perales.
2024a. \emph{Real Roots of Hypergeometric Polynomials via Finite Free
Convolution}. Version 3. \url{https://arxiv.org/abs/2309.10970v3}.

\bibitem[\citeproctext]{ref-martinez2024zeros}
Martínez-Finkelshtein, Andrei, Rafael Morales, and Daniel Perales.
2024b. \emph{Zeros of Generalized Hypergeometric Polynomials via Finite
Free Convolution. Applications to Multiple Orthogonality}. Version 2.
\url{https://arxiv.org/abs/2404.11479v2}.

\bibitem[\citeproctext]{ref-dlmf}
Olver, F. W. J., A. B. Olde Daalhuis, D. W. Lozier, et al., eds. 2026.
\emph{{NIST Digital Library of Mathematical Functions}}.
\href{https://dlmf.nist.gov/}{Https://dlmf.nist.gov/}, Release 1.2.7 of
2026-06-15. \url{https://dlmf.nist.gov/}.

\bibitem[\citeproctext]{ref-perry2016}
Perry, Kaitlyn Anne. 2016. {``Polydegree Properties of Polynomial
Automorphisms.''} Doctoral dissertation, The University of Alabama.
\url{https://ir.ua.edu/items/f71aa1b3-0dbc-42a2-9f59-3f01eb1a5529}.

\bibitem[\citeproctext]{ref-stacks00hp}
Stacks project authors, The. 2026a. \emph{The Stacks Project}.
\href{https://stacks.math.columbia.edu}{Https://stacks.math.columbia.edu}.
\url{https://stacks.math.columbia.edu/tag/00HP}.

\bibitem[\citeproctext]{ref-stacks00fv}
Stacks project authors, The. 2026b. \emph{The Stacks Project}.
\href{https://stacks.math.columbia.edu}{Https://stacks.math.columbia.edu}.
\url{https://stacks.math.columbia.edu/tag/00FV}.

\end{CSLReferences}

\end{document}
